\documentclass[11pt, letterpaper]{article}
\usepackage[margin=1in]{geometry}
\usepackage{enumitem}
\usepackage{parskip}
\usepackage{titlesec}
\usepackage[hidelinks]{hyperref}

\titleformat{\section}{\normalfont\bfseries}{}{0em}{}[\titlerule]
\titlespacing{\section}{0pt}{8pt}{4pt}

\begin{document}

\begin{center}
    {\Large \textbf{Sprout: A Light-Guided Growth Puzzle}}\\[2pt]
    {\small CIS 5680 Group Game Project 1 --- Spring 2026}
\end{center}

\vspace{-4pt}

\section*{High Concept}

You are a sentient plant rooted to the ground while a tiny traveler marches forward at a fixed pace
and cannot stop. Reposition light sources to steer your branches toward the light and away from
darkness, sculpting a living bridge across each gap before the traveler reaches the edge.

\section*{Genre}

2D side-scrolling light-guided growth puzzle.

\section*{Main Features}

\begin{itemize}[leftmargin=1.5em, itemsep=3pt]

    \item \textbf{You Are the Plant.} The player controls a living plant growing from a fixed root.
    The plant cannot be steered directly---instead, branches always grow toward the nearest light
    source. Moving, rotating, or toggling lights in the environment is the only way to guide where
    new growth goes. The plant's stem and branches form the sole bridge material available.

    \item \textbf{The Relentless Traveler.} A small creature walks automatically at a constant
    speed in whatever direction it currently faces. It cannot stop, jump, or be commanded. The
    player must coax the plant to grow a continuous path \textit{ahead} of the traveler before it
    reaches the next gap. Falling off any edge restarts the level.

    \item \textbf{Light-Guided Branching.} Each active light source pulls new growth toward it.
    Placing a lantern above a gap causes a branch to arc upward and across; moving the lantern
    left or right steers the arc. Blocking or extinguishing a light causes growth in that direction
    to stall and a different lit path to take over. Orchestrating multiple light sources at once
    shapes complex branching structures.

    \item \textbf{Fork and Redirect.} When two light sources compete, the plant forks---sending
    one branch toward each. Players exploit this to build multi-branch structures: one fork as a
    floor, another as a ramp, a third anchoring the span to a ledge. The puzzle is figuring out
    which light arrangement produces the fork geometry needed to span the level's obstacles.

    \item \textbf{Wither and Regrow.} Players can extinguish a light to stop and wither a branch
    that grew in the wrong direction, then relight it from a new angle to try again. This makes
    experimentation low-stakes and encourages iterative refinement of the plant's shape before
    the traveler arrives.

\end{itemize}

\section*{Player Challenge / Motivation}

The player thinks like a gardener and a puzzle solver: read the obstacle ahead, decide which light
arrangement will coax the plant into the right shape, and execute it before the traveler reaches
the gap. There is no direct building---every action is indirect, mediated through light. The reward
is watching a living bridge grow organically into exactly the shape the puzzle required.

\section*{Design Goals}

\begin{itemize}[leftmargin=1.5em, itemsep=2pt]
    \item \textbf{Indirect Control.} The player never touches the plant directly; all influence
    is through light placement. This creates a satisfying layer of indirection that makes
    solutions feel discovered rather than drawn.
    \item \textbf{Readable.} The direction of growth is always visually predictable---branches
    lean toward lit areas---so the player can plan ahead confidently without guessing.
    \item \textbf{Escalating.} Early levels use a single moveable lantern; later levels introduce
    fixed lights, lights on timers, and obstacles that block light mid-level, requiring the player
    to sequence light changes carefully.
\end{itemize}

\section*{Target Customer}

Fans of indirect-control and creative puzzle games who enjoy guiding systems rather than
micromanaging them directly. Players who appreciate a calm, iterative puzzle experience with a
clear visual language. Ages 12 and up; pick-up sessions of 5--15 minutes per level.

\section*{Unique Selling Points}

\begin{itemize}[leftmargin=1.5em, itemsep=2pt]
    \item The player steers a living plant entirely through light manipulation---no direct drawing
    or placement, just repositioning lanterns and watching nature respond.
    \item Forking branches emerge naturally from competing light sources, making multi-path
    structures an organic consequence of the core mechanic rather than a separate tool.
    \item A passively moving NPC traveler creates time pressure and emotional stakes without any
    combat or enemy AI.
\end{itemize}

\section*{Competition}

Indirect-control games such as \textit{Lemmings} guide a passive character through environmental
manipulation but do not use organic growth. Plant-growth games such as \textit{Flower} lack puzzle
structure and a dependent NPC. \textit{Sprout} is unique in using light placement as the sole
steering mechanism for a growing, branching structure that must reach a goal before a moving
character falls.

\section*{Target Hardware}

PC (Windows), implemented in Unreal Engine 5. Mouse-and-keyboard or gamepad supported.

\end{document}
